\documentclass{article}

\usepackage[utf8]{inputenc}
\usepackage[T1]{fontenc}
\usepackage{amsmath,amssymb,amsthm}
\usepackage{mathtools}
\usepackage{hyperref}
\usepackage{booktabs}

\makeatletter
\newcommand{\citep}[2][]{%
  \begingroup
  \def\@tmpkey{#2}%
  \ifx\@tmpkey\@empty\else
    (\ifx\@empty#1\@empty\else #1, \fi
     \@nameuse{bibkey@\@tmpkey})%
  \fi
  \endgroup
}
\expandafter\def\csname bibkey@lasser2026horizon\endcsname{Horizon~Aware}
\expandafter\def\csname bibkey@lasser2026exo\endcsname{Exogenous~Verification}
\expandafter\def\csname bibkey@lasser2026phase\endcsname{Phase~Redundancy}
\makeatother

\newtheorem{definition}{Definition}
\newtheorem{lemma}{Lemma}
\newtheorem{theorem}{Theorem}
\newtheorem{remark}{Remark}
\newtheorem{assumption}{Assumption}

\newcommand{\rext}{r_{\mathrm{ext}}}
\newcommand{\meff}{m_{\mathrm{eff}}}
\newcommand{\mindep}{m_{\mathrm{eff}}^{\mathrm{indep}}}

\title{Lemma 5a (Substrate floor): Full Proof Draft}
\author{Paper 10 working draft for codex review}
\date{}

\begin{document}

\maketitle

\section{Goal of this draft}

Promote Lemma~5a of Paper~10 \S4.7 from inline outline to a
self-contained formal proof for the appendix.

The current outline asserts that $\rext \geq r_*(m^*) > 0$ under
$I_6'$ ($\mindep \geq m^* \geq 3$) and substrate-distinctness, with
$r_*(m^*) \sim \rho_0 \binom{m^*}{2}$ for pairwise-cooperative-
dominated structures.  The outline references a ``coercivity-style
assumption'' that needs to be made explicit and a combinatorial
scaling argument that needs formal proof.

This draft formalizes the coercivity assumption (CN: cooperative
novelty rate floor), proves the pairwise-channel scaling under
independence, and discusses the sub-additive correlated-channel
caveat.

\section{Source-paper grounding}

\citep[Horizon~Aware]{lasser2026horizon}'s anti-monopolar property
establishes that under a discounted objective with $\rext > 0$,
domination is anti-maximizing.  $\rext$ is defined as the
cross-substrate cooperative novelty rate: the rate at which
cooperative capabilities form across distinct substrates that no
single substrate could produce alone.

\citep[Horizon~Aware]{lasser2026horizon} grounds $\rext > 0$ in
substrate physics (cross-substrate cooperatives are produced by
distinct-substrate participants whose joint output cannot be
replicated by any single substrate) but does not give a quantitative
floor.  Lemma~5a's contribution is to derive an explicit lower
bound on $\rext$ given the substrate count $\mindep$.

The bound is derived from two ingredients:
\begin{itemize}
\item A \emph{per-channel} coercivity assumption: each pair of
  distinct substrates contributes at least rate $\rho_0$ of
  cooperative novelty.
\item A \emph{combinatorial} structure: $\mindep$ pairwise-distinct
  substrates produce $\binom{\mindep}{2}$ pairwise cooperative
  channels.
\end{itemize}

Under joint failure-correlation independence (the substrate-
distinctness condition of $I_6'$, Definition~\ref{def:joint_indep}),
the rates add, giving the closed-form floor.

\section{Coercivity assumption}

\begin{assumption}[CN: Per-pair cooperative-novelty rate floor]
\label{ass:CN}
For every pair of distinct substrates $(s_i, s_j)$ in the
deployment's substrate population, the cross-substrate cooperative
novelty rate contributed by the pair $(s_i, s_j)$ alone is bounded
below by a positive constant $\rho_0 > 0$ depending on the
deployment class $\mathcal{D}$ but independent of $|P|$:
\[
  \rext^{(s_i, s_j)} \;\geq\; \rho_0 \;>\; 0.
\]
\end{assumption}

\textbf{Operational interpretation.}  $\rho_0$ measures the
\emph{baseline} per-pair cooperative-novelty rate: the rate at
which any pair of failure-correlation-independent substrates
generates cooperatives that neither substrate alone could produce.
$\rho_0$ depends on:
\begin{itemize}
\item Substrate-difference magnitude: substrates with more
  qualitatively different capability profiles produce more novelty
  per pair (e.g., human + AI cooperatives produce more novelty
  than two-similar-AI cooperatives, even if both pairs are formally
  $\mindep$-independent).
\item Cooperation infrastructure: deployments with
  governance-fork-supported cooperative protocols (\citep[Exogenous~Verification]{lasser2026exo})
  produce more cooperatives per substrate-pair than ad-hoc
  arrangements.
\item Capability-poset density: deployments with rich poset
  structure on $P^{\mathrm{act}}$ admit more cooperative
  combinations per substrate-pair than thin posets.
\end{itemize}

CN is structurally analogous to
\citep[Phase~Redundancy]{lasser2026phase}'s coercivity assumptions
(C1, C2), which assert positive correction and degradation rates
without deriving them from finer machinery.  Like those assumptions,
CN is taken as a deployment-class condition; deployments where
$\rho_0 \to 0$ for some substrate pair fall outside Lemma~5a's
scope.

\section{Statement of Lemma 5a}

\begin{lemma}[Substrate floor]
\label{lem:5a}
Under invariant $I_6'$ ($\mindep \geq m^* \geq 3$), the joint
failure-correlation independence condition
(Definition~\ref{def:joint_indep}, the substrate-distinctness
condition of (C2) of Theorem~\ref{thm:deployment_safety}), and
Assumption~\ref{ass:CN} (per-pair cooperative-novelty rate floor):
\[
  \rext \;\geq\; r_*(m^*) \;:=\; \rho_0 \cdot \binom{m^*}{2} \;>\; 0.
\]
For $m^* = 3$, $r_*(3) = 3\rho_0 > 0$.  For general $m^*$,
$r_*(m^*) \sim \rho_0 m^{*\,2}/2$ as $m^* \to \infty$.

The floor $r_*(m^*)$ is intensive in $|P|$ over the deployment
class $\mathcal{D}$: $\rho_0$ depends on $\mathcal{D}$'s operational
parameters but not on $|P|$.
\end{lemma}

\section{Proof}

\begin{proof}
We work over a deployment class $\mathcal{D}$ satisfying $I_6'$,
joint failure-correlation independence, and CN.

\emph{Step 1 (substrate counting).}  By $I_6'$, the deployment has
$\mindep \geq m^* \geq 3$ pairwise-distinct substrates in the joint
event-class failure-correlation-independence sense
(Definition~\ref{def:joint_indep}).  Let $\mathcal{S} = \{s_1,
\ldots, s_{\mindep}\}$ denote the maximal independent substrate set.

\emph{Step 2 (combinatorial pairwise channels).}  The number of
unordered pairs of distinct substrates in $\mathcal{S}$ is:
\[
  \binom{\mindep}{2} \;=\; \frac{\mindep \cdot (\mindep - 1)}{2}
  \;\geq\; \binom{m^*}{2}
  \;=\; \frac{m^* (m^* - 1)}{2}.
\]
Each pair $(s_i, s_j)$ with $i < j$ defines a pairwise cooperative
channel: the set of cooperative capabilities whose participants
include at least one agent from $s_i$ and at least one agent from
$s_j$.

\emph{Step 3 (per-pair rate).}  By Assumption~CN, each pair
$(s_i, s_j)$ contributes a rate $\rext^{(s_i, s_j)} \geq \rho_0 >
0$ of cooperative novelty.

\emph{Step 4 (rate additivity under joint independence).}  Under
joint failure-correlation independence
(Definition~\ref{def:joint_indep}), the per-pair rates are
\emph{additive}: the cross-substrate cooperative novelty rate
across the substrate population satisfies:
\[
  \rext \;\geq\; \sum_{1 \leq i < j \leq \mindep}
  \rext^{(s_i, s_j)}.
\]
The additivity follows from the fact that joint independence
implies the per-pair cooperative outputs are statistically
independent: a cooperative produced by pair $(s_i, s_j)$ is not
counted in the rate of pair $(s_k, s_l)$ unless $\{i, j\} \cap \{k,
l\} \neq \emptyset$, in which case it is a higher-order cooperative
not in the pairwise count.  Higher-order cooperatives contribute
\emph{additionally} to $\rext$, so the pairwise sum is a lower
bound.

\emph{Step 5 (closed-form floor).}  Combining Steps 1--4:
\[
  \rext \;\geq\; \sum_{1 \leq i < j \leq \mindep} \rext^{(s_i, s_j)}
  \;\geq\; \rho_0 \cdot \binom{\mindep}{2}
  \;\geq\; \rho_0 \cdot \binom{m^*}{2}
  \;=:\; r_*(m^*).
\]
For $m^* = 3$:
\[
  r_*(3) \;=\; \rho_0 \cdot \binom{3}{2} \;=\; 3 \rho_0 \;>\; 0.
\]

\emph{Step 6 (intensivity).}  $\rho_0$ is a deployment-class
constant (CN), independent of $|P|$.  $m^*$ is the threshold from
$I_6'$, also a deployment-class constant.  Therefore $r_*(m^*) =
\rho_0 \binom{m^*}{2}$ is independent of $|P|$ over $\mathcal{D}$.
$\Box$
\end{proof}

\section{Caveats}

\paragraph{Sub-additive correlated channels.}
The proof's Step~4 uses joint failure-correlation independence to
derive rate additivity.  When pairs of substrates have
\emph{correlated} cooperative-channel structure (e.g., two
substrates contribute redundant cooperatives that one of them
already supplied), the per-pair rates are sub-additive:
\[
  \rext \;<\; \sum \rext^{(s_i, s_j)} \quad \text{(correlated case)}.
\]
The lemma's bound $r_*(m^*) = \rho_0 \binom{m^*}{2}$ is then
\emph{anti-conservative} (over-estimates $\rext$).  Under joint
failure-correlation independence (which is precisely the (C2)
requirement of Theorem~\ref{thm:deployment_safety} and the
$\mindep$ formulation of $I_6'$), correlated-channel sub-additivity
is excluded.

The lemma is therefore strictly conditional on joint independence;
deployments where the substrate-distinctness audit (Audit~1 of
\S\ref{sec:tooling_substrate_audits}) fails to verify joint
independence cannot invoke Lemma~5a directly.

\paragraph{Higher-order cooperatives.}
The proof bounds $\rext$ from below using only \emph{pairwise}
cooperative channels.  Higher-order cooperatives (those requiring
participants from three or more distinct substrates) contribute
additional rate, increasing $\rext$ above $r_*(m^*)$.  For canonical
tripartite deployments (Human + AI + Formal-Operational), the
three-way cooperative ``verified workflow'' adds an additional
$\rho_0^{(3)}$ rate (the three-way coercivity constant), yielding
$\rext \geq 3 \rho_0 + \rho_0^{(3)} > 3\rho_0$.

The lemma's pairwise floor is therefore a conservative lower bound;
deployments may calibrate higher floors using their specific
cooperative structure.  The pairwise scaling $\binom{m^*}{2}$ is
the universal lower bound applicable to any deployment satisfying
$I_6'$.

\paragraph{Per-pair heterogeneity.}
CN states a uniform floor $\rho_0$ across pairs.  In practice,
different substrate pairs produce different cooperative-novelty
rates: e.g., Human-AI pair may produce $\rho_{\mathrm{HA}}$,
Human-Formal pair $\rho_{\mathrm{HF}}$, AI-Formal pair
$\rho_{\mathrm{AF}}$, with $\rho_0 = \min\{\rho_{\mathrm{HA}},
\rho_{\mathrm{HF}}, \rho_{\mathrm{AF}}\}$.  The proof uses the
minimum, giving the conservative bound; deployments may use
heterogeneous per-pair rates for tighter floors.

\section{Connection to Theorem 1 Layer 1}

Lemma~5a's role: provide the substrate-floor input that makes
Lemma~5c's safe-region inequality non-trivial.

Lemma~5c asserts:
\[
  \Delta r_K \;<\; \rext + (1 - \gamma)(\Ddiv - \Delta_0)
\]
implies diversity dominates.  For this safe region to be
non-trivial, $\rext$ must be bounded below by a deployment-
verifiable constant.  Lemma~5a provides:
\[
  \rext \;\geq\; \rho_0 \binom{m^*}{2} \;>\; 0,
\]
making the safe region non-trivial whenever $\rho_0$ is
operationally calibrated.

The connection is direct: Lemma~5c invokes Lemma~5a in its proof
(Step~1 of the Lemma~5c appendix proof) to assert $\rext \geq
r_*(m^*) > 0$ before applying the strategy-comparison machinery.

\section{What this proof does and does not establish}

\textbf{Establishes:}
\begin{itemize}
\item A formal coercivity assumption (CN: per-pair cooperative-
  novelty rate floor $\rho_0 > 0$) consistent with the structural
  pattern of \citep[Phase~Redundancy]{lasser2026phase}'s C1, C2.
\item A combinatorial-additivity proof of the pairwise scaling
  $\binom{m^*}{2}$ under joint failure-correlation independence.
\item An explicit closed-form floor $r_*(m^*) = \rho_0 \binom{m^*}{2}$
  intensive in $|P|$.
\item Three caveats: sub-additive correlated-channel case,
  higher-order cooperatives, per-pair heterogeneity.
\end{itemize}

\textbf{Does not establish (deferred):}
\begin{itemize}
\item Tightness of $r_*(m^*)$.  The pairwise-only floor is
  conservative; higher-order cooperatives contribute additional
  rate.
\item Existence of deployment configurations satisfying CN with
  large $\rho_0$.  CN is a condition on the deployment's substrate
  structure; verification is empirical (calibration via the
  deployment-tooling specification).
\item Numerical value of $\rho_0$.  Calibrated deployment-by-
  deployment via the per-pair cooperative-novelty rate-estimation
  procedure.
\end{itemize}

\section{Specific verification questions for codex review}

\begin{enumerate}
\item Is Assumption CN (per-pair cooperative-novelty rate floor)
  the right formalization?  The assumption asserts a uniform $\rho_0$
  across pairs.  Should it admit per-pair heterogeneity by default,
  with $\rho_0 = \min_{i,j} \rho^{(s_i, s_j)}$ derived rather than
  asserted?

\item Is Step~4 (rate additivity under joint independence) rigorous?
  The proof argues per-pair cooperatives are statistically
  independent under
  Definition~\ref{def:joint_indep}.  Is this argument
  watertight, or does it require additional structural conditions
  (e.g., that cooperative production from one pair does not
  consume capabilities that would otherwise contribute to another
  pair's cooperatives)?

\item Is the higher-order cooperative caveat correctly handled?
  The proof bounds $\rext$ using only pairwise channels, with
  higher-order channels contributing additionally.  Is this the
  right way to handle them, or should they be incorporated into
  CN directly (e.g., $\rho_0^{(k)}$ for each $k$-way cooperative
  with $\rext \geq \sum_k \rho_0^{(k)} \binom{m^*}{k}$)?

\item Is the $r_*(m^*) = \rho_0 \binom{m^*}{2}$ formula consistent
  with \citep[Horizon~Aware]{lasser2026horizon}'s anti-monopolar
  framework?  In particular, does Paper~3 give any quantitative
  bound on $\rext$ that this lemma is in tension with, or is the
  lemma a strict refinement?

\item Should CN be promoted to a theorem-level condition (e.g.,
  (C10)) of Theorem~\ref{thm:deployment_safety}, analogous to
  the (C6)-(C9) conditions added in Lemma~1 and Lemma~2's
  integration?  Or can it remain a Lemma~5a deployment-class
  assumption without theorem-level surfacing?
\end{enumerate}

\end{document}
